\documentclass{article}
\usepackage[nonanonymous]{neurips_2026}
\usepackage[utf8]{inputenc}
\usepackage[T1]{fontenc}
\usepackage{microtype}
\usepackage{booktabs}
\usepackage{graphicx}
\usepackage{tabularx}
\usepackage{array}
\usepackage{hyperref}
\hypersetup{colorlinks=true,linkcolor=blue,urlcolor=blue}
\newcommand{\figdir}{figures}
\title{Supplementary Material: Beyond Replay RMSE}
\author{Anonymous Authors}
\begin{document}
\maketitle

\section{Condition-level observer selection}
Table~\ref{tab:conditions} reports the closed-loop oracle and the observers chosen by replay position RMSE and LCSE. Regret is the increase in closed-loop cross-track RMSE relative to the oracle. The corresponding regret matrix is shown in Figure~\ref{fig:regret}.

\begin{table}[h]
\centering
\caption{Condition-level selection. Regret is in centimetres.}
\label{tab:conditions}
\scriptsize
\begin{tabular}{lccrcr}
\toprule
Condition & CL oracle & Pose choice & Pose regret & LCSE choice & LCSE regret\\
\midrule
Blackout 4 & GRU-EKF & GRU-EKF & 0.00 & GRU-EKF & 0.00 \\
Blackout 10 & GRU-EKF & GRU-EKF & 0.00 & GRU-EKF & 0.00 \\
Blackout 16 & GRU-EKF & GRU-EKF & 0.00 & GRU-EKF & 0.00 \\
Blackout 22 & GRU-EKF & GRU-EKF & 0.00 & GRU-EKF & 0.00 \\
Blackout 28 & EKF & SSM-EKF & 2.35 & EKF & 0.00 \\
Blackout 34 & EKF & SSM-EKF & 2.93 & EKF & 0.00 \\
Range noise 0.05 & EKF & EKF & 0.00 & GRU-EKF & 1.39 \\
Range noise 0.12 & GRU-EKF & GRU-EKF & 0.00 & GRU-EKF & 0.00 \\
Range noise 0.2 & GRU-EKF & GRU-EKF & 0.00 & GRU-EKF & 0.00 \\
Range noise 0.3 & GRU-EKF & GRU-EKF & 0.00 & GRU-EKF & 0.00 \\
Range noise 0.4 & GRU-EKF & SSM-EKF & 2.07 & GRU-EKF & 0.00 \\
Range noise 0.5 & GRU-EKF & SSM-EKF & 3.47 & GRU-EKF & 0.00 \\
Landmarks 2 & SSM-EKF & SSM-EKF & 0.00 & GRU-EKF & 9.40 \\
Landmarks 3 & GRU-EKF & GRU-EKF & 0.00 & GRU-EKF & 0.00 \\
Landmarks 4 & GRU-EKF & GRU-EKF & 0.00 & GRU-EKF & 0.00 \\
Landmarks 5 & GRU-EKF & GRU-EKF & 0.00 & GRU-EKF & 0.00 \\
Landmarks 6 & GRU-EKF & GRU-EKF & 0.00 & GRU-EKF & 0.00 \\
Landmarks 8 & SSM-EKF & GRU-EKF & 1.40 & GRU-EKF & 1.40 \\
Gyro bias 0.0 & GRU-EKF & GRU-EKF & 0.00 & GRU-EKF & 0.00 \\
Gyro bias 0.05 & GRU-EKF & GRU-EKF & 0.00 & GRU-EKF & 0.00 \\
Gyro bias 0.1 & GRU-EKF & GRU-EKF & 0.00 & GRU-EKF & 0.00 \\
Gyro bias 0.15 & GRU-EKF & GRU-EKF & 0.00 & GRU-EKF & 0.00 \\
Gyro bias 0.2 & GRU-EKF & GRU-EKF & 0.00 & GRU-EKF & 0.00 \\
Gyro bias 0.25 & GRU-EKF & GRU-EKF & 0.00 & GRU-EKF & 0.00 \\
\bottomrule
\end{tabular}
\end{table}

\begin{figure}[h]
\centering
\includegraphics[width=.98\linewidth]{\figdir/fig_proxy_regret_matrix.png}
\caption{Closed-loop selection regret for each offline proxy and condition. B denotes blackout length, R range noise, L landmark count, and G gyro bias.}
\label{fig:regret}
\end{figure}

\section{Matched ablations and training initialization}
The paired intervals in Table~\ref{tab:abl_full} are computed over matched environment seeds and report the difference relative to the selective+mask model. Negative values favor the ablation. Figure~\ref{fig:abl_seed} visualizes the same-anchor results and the three independently trained SSM-EKF checkpoints.

\begin{table}[h]
\centering
\caption{Full matched EKF-anchor ablation results. Values are mean $\pm$ 95\% CI half-width; difference intervals are paired ablation-minus-selective intervals.}
\label{tab:abl_full}
\scriptsize
\begin{tabular}{lccccc}
\toprule
Condition & Selective & Non-selective & Paired $\Delta$ CI & No global mask & Paired $\Delta$ CI\\
\midrule
Nominal & 0.199$\pm$0.034 & 0.203$\pm$0.034 & [-0.051,0.059] & 0.195$\pm$0.028 & [-0.044,0.031] \\
Long blackout & 0.336$\pm$0.070 & 0.235$\pm$0.027 & [-0.182,-0.020] & 0.308$\pm$0.060 & [-0.102,0.040] \\
High range noise & 0.307$\pm$0.090 & 0.322$\pm$0.043 & [-0.093,0.103] & 0.294$\pm$0.044 & [-0.110,0.067] \\
Two landmarks & 0.289$\pm$0.050 & 0.235$\pm$0.024 & [-0.103,-0.007] & 0.274$\pm$0.038 & [-0.050,0.018] \\
High gyro bias & 0.314$\pm$0.065 & 0.393$\pm$0.039 & [0.011,0.145] & 0.312$\pm$0.039 & [-0.055,0.051] \\
\bottomrule
\end{tabular}
\end{table}

\begin{figure}[h]
\centering
\begin{minipage}[t]{.49\linewidth}\centering
\includegraphics[width=\linewidth]{\figdir/fig_ablation_compact.png}
\end{minipage}\hfill
\begin{minipage}[t]{.49\linewidth}\centering
\includegraphics[width=\linewidth]{\figdir/fig_seed_compact.png}
\end{minipage}
\caption{Same-anchor ablation (left) and training-initialization sensitivity (right).}
\label{fig:abl_seed}
\end{figure}

\begin{table}[h]
\centering
\caption{Across-training-seed SSM-EKF summary.}
\small
\begin{tabular}{lccc}
\toprule
Condition & Minimum mean & Maximum mean & Across-seed SD\\
\midrule
Nominal & 0.189 & 0.199 & 0.006\\
Long blackout & 0.231 & 0.336 & 0.060\\
High range noise & 0.249 & 0.307 & 0.033\\
Two landmarks & 0.226 & 0.289 & 0.032\\
High gyro bias & 0.293 & 0.314 & 0.011\\
\bottomrule
\end{tabular}
\end{table}

\section{Controller sensitivity}
The five-seed Kanayama sensitivity is not used as the primary statistical result. It tests whether conclusions transfer when the same observers, plant, actuator limits, conditions, and random seeds are paired with a different tracking law.
\begin{figure}[h]
\centering
\includegraphics[width=.68\linewidth]{\figdir/fig_controller_sensitivity.png}
\caption{Replay-RMSE selection failures for Kanayama and the fraction of conditions in which the pure-pursuit and Kanayama closed-loop oracles differ.}
\end{figure}
\begin{table}[h]
\centering
\small
\begin{tabular}{lc}
\toprule
Statistic & Result\\
\midrule
Replay-RMSE failures for Kanayama & 7/24 (29.2\%)\\
Mean / maximum Kanayama regret & 0.0142 / 0.0923 m\\
Closed-loop oracle changes & 8/24 (33.3\%)\\
Oracle-change bootstrap CI & 0.208--0.625\\
\bottomrule
\end{tabular}
\end{table}

\section{Full MRCLAM EKF grid and calibrated replay}
\begin{table}[h]
\centering
\caption{Physical-log localization on the 15-landmark MRCLAM segment.}
\label{tab:mr_grid}
\scriptsize
\begin{tabular}{lrrrr}
\toprule
Observer / covariance scale & Pos. RMSE & Head. RMSE & Command diff. & LCSE\\
\midrule
DR & 3.392 & 1.662 & 0.586 & 62.91 \\
EKF q0.25 r0.25 & 1.599 & 0.864 & 0.501 & 38.13 \\
EKF q0.25 r1 & 1.511 & 0.879 & 0.508 & 32.10 \\
EKF q0.25 r4 & 1.357 & 0.850 & 0.525 & 4.22 \\
EKF q1 r0.25 & 1.703 & 0.879 & 0.489 & 39.58 \\
EKF q1 r1 & 1.599 & 0.864 & 0.501 & 38.13 \\
EKF q1 r4 & 1.507 & 0.878 & 0.507 & 32.10 \\
EKF q4 r0.25 & 1.804 & 0.881 & 0.489 & 36.31 \\
EKF q4 r1 & 1.702 & 0.879 & 0.489 & 39.58 \\
EKF q4 r4 & 1.597 & 0.863 & 0.500 & 38.13 \\
\bottomrule
\end{tabular}
\end{table}
The best position result is obtained with process scale $q=0.25$ and measurement scale $r=4$. In the separate real-noise-calibrated simulated replay, dead reckoning is the closed-loop oracle. Pose RMSE selects it with zero regret, whereas command disagreement chooses EKF $q=1,r=4$ with 0.120 m regret and LCSE chooses EKF $q=0.25,r=4$ with 0.203 m regret.

\section{Reproducibility map}
\begin{table}[h]
\centering
\caption{Primary commands in the released package.}
\small
\begin{tabularx}{\linewidth}{lX}
\toprule
Task & Command\\
\midrule
Train checkpoints & \texttt{python scripts/run\_train.py --backend torch --force}\\
Synthetic evaluation & \texttt{python scripts/run\_eval.py}\\
Replay/closed-loop proxy audit & \texttt{python scripts/run\_proxy\_analysis.py}\\
Matched EKF ablations & \texttt{python scripts/run\_matched\_ablation.py}\\
Training-seed study & \texttt{python scripts/run\_training\_seed\_study.py}\\
Second controller & \texttt{python scripts/run\_second\_controller.py --seeds 5}\\
Download MRCLAM & \texttt{python scripts/download\_mrclam.py}\\
Physical-log analysis & \texttt{python scripts/run\_mrclam.py}\\
Publication figures & \texttt{python scripts/make\_enhanced\_figures.py}\\
\bottomrule
\end{tabularx}
\end{table}

\end{document}
